\section{Prototypische Realisierung}
\label{sec:implementation}

Das in Kapitel~\ref{sec:architecture} entwickelte Architekturkonzept wird
zur technischen Konkretisierung und Untersuchung prototypisch realisiert.
Der im Entwicklungsprojekt als \emph{Turing Generator} bezeichnete Prototyp
setzt die zuvor beschriebenen Verarbeitungsschritte und menschlichen
Entscheidungen technisch um. Im weiteren Verlauf wird er daher verkürzt als
\emph{Prototyp} bezeichnet.

Das in CATIA Magic entwickelte SysML-v2-Modell beschreibt hierzu die
Anforderungen an das System, die dafür erforderlichen Funktionen und deren
Zuordnung zu logischen Systemelementen. Die Softwareimplementierung setzt
diese für die Untersuchung erforderlichen Funktionen ausführbar um.

Auf hoher Ebene verarbeitet der Prototyp zunächst Engineering-Quellen, also
die als Eingang verwendeten Engineering-Dokumente und -daten, und bereitet
deren Inhalte für die fachliche Prüfung auf. Fachlich geprüfte und
autorisierte Engineering-Information wird anschließend in eine interne
Repräsentation des entstehenden Modells überführt und schließlich in
SysML~v2 dargestellt und validiert. Menschliche Entscheidungen steuern dabei
die fachliche Autorisierung, die Einordnung in das Modell sowie die
abschließende Modellprüfung und Freigabe. Bei der Verarbeitung mehrerer
Quellen werden deren Informationen zusätzlich auf Projektebene gemeinsam
betrachtet, ohne ihre Herkunft und frühere Entscheidungen zu verlieren.

Die folgenden Abschnitte erläutern zunächst die Systemarchitektur und danach
die einzelnen Verarbeitungsschritte bis zum erzeugten SysML-v2-Modell.
Abschließend werden die Verarbeitung mehrerer Quellen sowie die dauerhafte
Speicherung und Nachverfolgbarkeit der entstehenden Zustände beschrieben.

\subsection{Systemarchitektur des Prototyps}
\label{sec:implementation_system_architecture}

Ein früher Strukturierungsansatz unterschied drei Bereiche, die im
Architekturmodell als \emph{Data Layer}, \emph{Process Layer} und
\emph{Knowledge Layer} bezeichnet wurden. Diese Dreiteilung stellt keine
lineare Verarbeitungskette dar. Sie diente zunächst dazu, unterschiedliche
Aufgaben des Systems zu gruppieren. Abbildung~\ref{fig:implementation_three_layer_mapping}
zeigt, wie diese frühe Struktur in die spätere logische Systemarchitektur
überführt wurde.

\begin{figure}[H]
    \centering
    \includegraphics[
        width=\textwidth,
        height=0.76\textheight,
        keepaspectratio
    ]{Graphiken/5.2a_Three_Layer_Architecture_Mapping.pdf}
    \caption[Überführung der ursprünglichen Drei-Schichten-Struktur]{
        Überführung der ursprünglichen Drei-Schichten-Struktur in das
        logische Systemmodell. Der \emph{Data Layer} umfasst die Verarbeitung
        und Nachverfolgbarkeit von Engineering-Information sowie den Projekt-
        und Quellenkontext. Der \emph{Process Layer} umfasst die Steuerung des
        Verarbeitungsablaufs, menschliche Entscheidungen, Modellbildung sowie
        SysML-v2-Erzeugung und Validierung. Der \emph{Knowledge Layer} stellt
        unterstützende Informationen für mehrere dieser Funktionen bereit und
        wirkt daher querschnittlich. Eigene Darstellung auf Basis des
        entwickelten SysML-v2-Architekturmodells.}
    \label{fig:implementation_three_layer_mapping}
\end{figure}

Die querschnittliche Ausprägung des \emph{Knowledge Layer} greift die in
Abschnitt~\ref{sec:architecture_information_authority} entwickelte Trennung
zwischen den fachlich maßgeblichen Engineering-Informationen und
unterstützenden Kontextinformationen auf.

Auf dieser Grundlage wurde die Systemarchitektur auf logischer Ebene weiter
konkretisiert. Abbildung~\ref{fig:implementation_logical_architecture} zeigt,
welche Systemkomponenten dafür vorgesehen sind, wie sie Informationen
untereinander austauschen und welche Beteiligten oder Informationsquellen
außerhalb des eigentlichen Prototyps angebunden sind.

\begin{figure}[H]
    \centering
    \includegraphics[
        width=\textwidth,
        height=0.82\textheight,
        keepaspectratio
    ]{Graphiken/5.2b_Logical_Architecture.pdf}
    \caption[Logische Systemarchitektur des Prototyps]{
            Logische Systemarchitektur des Prototyps. Die internen Komponenten
            trennen Benutzerinteraktion, Projekt- und Quellenverwaltung,
            Ablaufsteuerung, Verarbeitung von Engineering-Information,
            menschliche Prüfung und Autorisierung, Herkunft und Nachverfolgbarkeit,
            Modellbildung, SysML-v2-Erzeugung sowie Validierung. Das Sprachmodell,
            der menschliche Reviewer und unterstützende Referenzinformationen liegen
            außerhalb des eigentlichen Prototyps und sind über definierte
            Schnittstellen angebunden. Eigene Darstellung auf Basis des entwickelten
            SysML-v2-Architekturmodells.}
    \label{fig:implementation_logical_architecture}
\end{figure}

Die logische Zerlegung überführt die in
Kapitel~\ref{sec:architecture} entwickelten Architekturmechanismen in
konkrete Systemkomponenten. Die Verarbeitung von Engineering-Information,
menschliche Autorisierung, Modellbildung, SysML-v2-Erzeugung und Validierung
bleiben dabei strukturell getrennt. Das Sprachmodell und die menschliche
Prüfung sind ebenso wie die semantischen und modellbezogenen
Referenzinformationen über definierte Schnittstellen in diese Struktur
eingebunden.

Für die detailliertere Modellierung wurde die Systemarchitektur anschließend
weiter untergliedert. Dabei wurden Anforderungen, erforderliche Funktionen
und die zugehörigen logischen Systemelemente miteinander verknüpft
(R/F/L-Modellierung). Hierfür ist das System im Modell in neun Subsysteme,
also neun funktional abgegrenzte Teilbereiche, gegliedert. Die im Modell
verwendeten Bezeichnungen lauten:

\begin{enumerate}
    \item \emph{User Interaction and Guided Workflow},
    \item \emph{Project and Source Context Management},
    \item \emph{Processing Orchestration and State Control},
    \item \emph{Engineering Information Processing},
    \item \emph{Human Review and Engineering Authority},
    \item \emph{Evidence Coverage and Traceability},
    \item \emph{Architecture Derivation and Model Assembly},
    \item \emph{SysML v2 Artifact Generation},
    \item \emph{Generated Artifact Validation and Publication}.
\end{enumerate}

Diese Bezeichnungen dienen der Strukturierung des SysML-v2-Modells. Inhaltlich
decken die Subsysteme die zuvor beschriebenen Funktionen von der
Quellenverarbeitung über menschliche Autorisierung und Modellbildung bis zur
Erzeugung, Validierung und Freigabe des SysML-v2-Artefakts ab.

Diese Untergliederung verfeinert die logische Systemarchitektur für die
nachfolgende Modellierung und technische Umsetzung. Die folgenden Abschnitte
geben nicht die vollständige Modellhierarchie wieder, sondern konzentrieren
sich auf den tatsächlichen Verarbeitungsablauf des Prototyps.

\subsection{Durchgängige Realisierung des Informationsflusses}
\label{sec:implementation_e2e_flow}

Die in Abschnitt~\ref{sec:implementation_system_architecture}
beschriebenen Systemkomponenten wirken gemeinsam in einem durchgängigen
Verarbeitungsablauf. Dieser beginnt mit der Registrierung einer
Engineering-Quelle und endet mit einem validierten und explizit
freigegebenen SysML-v2-Artefakt.

Zwischen Eingabe und Ergebnis entstehen mehrere klar voneinander getrennte
Zwischenzustände. Der Inhalt einer Quelle wird zunächst aufbereitet und auf
relevante Engineering-Information untersucht. Diese Information wird
anschließend fachlich geprüft und autorisiert. Erst danach wird entschieden,
wie sie in das interne Modell eingeordnet werden kann. Auf dieser Grundlage
entsteht schrittweise ein interner Modellzustand, aus dem schließlich die
SysML-v2-Repräsentation erzeugt, validiert und abschließend menschlich
geprüft und freigegeben wird. Der Prototyp wandelt Eingangsdokumente damit
nicht unmittelbar in SysML-v2-Code um.

Abbildung~\ref{fig:implementation_e2e_flow} zeigt diesen Ablauf auf
Systemebene. Die im CATIA-Modell als
\emph{SFB\_004 Engineering Data Transformation Flow} bezeichnete Sicht
fasst die wesentlichen Verarbeitungsschritte und ihre Informationsübergänge
zusammen. Die englischen Bezeichnungen in der Abbildung entsprechen dabei
den Namen der Funktionen und Informationsobjekte im SysML-v2-Modell.

\begin{figure}[H]
    \centering
    \includegraphics[
        width=\textwidth,
        height=0.82\textheight,
        keepaspectratio
    ]{Graphiken/5.3_End_to_End_Engineering_Information_Flow.pdf}
    \caption[Durchgängiger Informationsfluss des Prototyps]{
        Durchgängiger Informationsfluss von der registrierten
        Engineering-Quelle bis zum freigegebenen SysML-v2-Artefakt.
        Nach der Aufbereitung und fachlichen Prüfung der Quellinformation
        wird deren Eignung für die projektweite Verarbeitung bewertet.
        Anschließend werden daraus Modellinhalte abgeleitet, menschlich
        autorisiert und in einem internen Modellzustand zusammengeführt.
        Dieser bildet die Grundlage für die SysML-v2-Erzeugung, deren
        Validierung sowie die abschließende menschliche Modellprüfung und
        Freigabe. Sprachmodell und menschlicher Reviewer sind als externe
        Beteiligte dargestellt. Eigene Darstellung auf Basis des
        entwickelten SysML-v2-Architekturmodells.}
    \label{fig:implementation_e2e_flow}
\end{figure}

Der Ablauf beginnt mit der Registrierung einer Engineering-Quelle. Dabei
wird festgehalten, welchem Projekt die Quelle zugeordnet ist, um welche
konkrete Quelle es sich handelt und unter welchen Bedingungen sie verarbeitet
wird. Diese Angaben bilden den Projekt-, Quellen- und Verarbeitungskontext.
Sie stellen sicher, dass später erzeugte Informationen und Modellartefakte
auf die jeweilige Quelle und den zugehörigen Verarbeitungslauf
zurückgeführt werden können. Die technische Speicherung dieser Informationen und
ihrer Verknüpfungen wird in
Abschnitt~\ref{sec:implementation_persistence_limits} beschrieben.

Anschließend wird der Inhalt der Quelle verarbeitet und in eine Form
überführt, in der relevante Engineering-Information gezielt geprüft werden
kann. Die Verarbeitung selbst autorisiert noch keine fachlichen Aussagen.
Diese Autorisierung erfolgt erst durch eine menschliche Prüfung
(\emph{Human Engineering Authority}). Dadurch bleibt nachvollziehbar,
welche Information lediglich aus der Quelle abgeleitet wurde und welche
Information anschließend fachlich bestätigt wurde.

Bevor autorisierte Information gemeinsam mit weiteren Quellen auf
Projektebene verarbeitet wird, wird ihre Eignung für den betrachteten
Projektkontext geprüft. Dieser Schritt wird im Modell als
\emph{Project Fit} bezeichnet. Er beantwortet damit nicht die Frage, ob
eine Aussage fachlich richtig ist, sondern ob sie für das betrachtete
Projekt in die weitere Verarbeitung einbezogen werden soll.

Aus den zugelassenen Engineering-Informationen werden anschließend mögliche
Modellelemente und Beziehungen abgeleitet. Diese Vorschläge werden nicht
unmittelbar in das Modell übernommen. Die menschliche
\emph{Model Authority} entscheidet zunächst über ihre strukturelle
Einordnung. Erst auf Grundlage dieser Entscheidung wird der interne
Modellzustand aufgebaut.

Dieser interne Modellzustand wird anschließend für die Zielrepräsentation
aufbereitet und deterministisch in SysML~v2 überführt. Das erzeugte
SysML-v2-Artefakt wird validiert und danach einer abschließenden
menschlichen Modellprüfung unterzogen. Erst nach dieser Freigabe gilt der
Verarbeitungspfad als abgeschlossen.

Die folgenden Abschnitte erläutern die einzelnen Teile dieses Ablaufs
detaillierter. Abschnitt~\ref{sec:implementation_information_processing}
beschreibt die Verarbeitung der Engineering-Quellen bis zur fachlichen
Prüfung. Abschnitt~\ref{sec:implementation_authority} behandelt die
menschlichen Entscheidungen über fachliche Information und Modellbildung.
Abschnitt~\ref{sec:implementation_model_realization} beschreibt den Aufbau
des internen Modells sowie dessen Überführung in SysML~v2. Die Verarbeitung
mehrerer Engineering-Quellen wird in
Abschnitt~\ref{sec:implementation_multisource} ergänzt.

Der dargestellte Ablauf enthält außerdem kontrollierte Rücksprünge und
Abbruchbedingungen. Fehlen beispielsweise erforderliche Informationen oder
eine notwendige menschliche Entscheidung, wird der nachfolgende
Verarbeitungsschritt nicht ausgeführt. Damit wird das in
Abschnitt~\ref{sec:architecture_controlled_transformation} eingeführte
\emph{fail-closed}-Prinzip umgesetzt.

Informationen über Herkunft, Verarbeitung und getroffene Entscheidungen
werden während des gesamten Ablaufs mitgeführt und dauerhaft gespeichert.
Dadurch bleiben die in Kapitel~\ref{sec:architecture} beschriebenen
Provenienz- und Nachverfolgbarkeitsbeziehungen auch in der prototypischen
Realisierung erhalten.

\paragraph{Technische Verträge zwischen Verarbeitungsschritten}

Die Übergänge zwischen den einzelnen Verarbeitungsschritten werden in der
prototypischen Realisierung durch explizite technische Verträge
konkretisiert. Ein solcher Vertrag legt fest, welche Eingangszustände ein
Verarbeitungsschritt verwenden darf, welche Identitäten und Referenzen
erhalten bleiben müssen, welche Bedingungen für eine gültige Verarbeitung
gelten und welche Ergebniszustände für die nachfolgende Verarbeitung
bereitgestellt werden können.

Die Verträge übertragen damit die im Architekturkonzept definierten
Informations-, Provenienz- und Autorisierungsgrenzen in überprüfbare
Bedingungen der Implementierung. Sie verhindern beispielsweise, dass
Persona-Interpretationen auf unterschiedlichen fachlichen
Referenzgegenständen beruhen, dass menschliche Entscheidungen von dem
tatsächlich geprüften Informationszustand gelöst werden oder dass ein
Modellzustand ohne die erforderliche fachliche und strukturelle
Autorisierung weiterverarbeitet wird.

Solche Verträge werden im Prototyp unter anderem für die Bildung der
fachlichen Referenzgegenstände (\emph{Engineering Subjects}), die
Persona-basierte Interpretation, den Vergleich der
Interpretationsergebnisse, die menschliche fachliche Autorisierung und die
Modellplatzierung verwendet. In der Implementierung werden diese
Festlegungen als \emph{Contracts} bezeichnet.

Die technischen Verträge bilden zugleich eine Grundlage der automatisierten
Verifikation. Ihre regelbasiert überprüfbaren Bedingungen werden in
Abschnitt~\ref{sec:validation_software} als Bestandteil der Software- und
Vertragsverifikation untersucht. Eine Übersicht der für die prototypische
Realisierung wesentlichen Verträge ist in
Anhang~\ref{app:implementation_contracts} dokumentiert.

\subsection{Quellenbezogene Verarbeitung von Engineering-Information}
\label{sec:implementation_information_processing}

Der erste Teil des in
Abschnitt~\ref{sec:implementation_e2e_flow} beschriebenen
Verarbeitungsablaufs bereitet den Inhalt einer einzelnen Engineering-Quelle
für die nachfolgende fachliche Prüfung auf. Im CATIA-Modell ist dieser
Bereich als \emph{SFB\_005 Engineering Information Processing Flow}
bezeichnet. Er reicht von der registrierten Quelle bis zu den Informationen,
die dem menschlichen Reviewer für seine Entscheidung bereitgestellt werden.
Abbildung~\ref{fig:implementation_information_processing} zeigt die dafür
realisierte Verarbeitungskette.

\begin{figure}[H]
    \centering
    \includegraphics[
        width=\textwidth,
        height=0.84\textheight,
        keepaspectratio
    ]{Graphiken/5.4_Engineering_Information_Processing_Flow.pdf}
    \caption[Quellenbezogene Verarbeitung von Engineering-Information]{
        Quellenbezogene Verarbeitung von Engineering-Information bis zur
        menschlichen fachlichen Prüfung. Der Inhalt einer registrierten
        Quelle wird zunächst technisch aufbereitet und in eindeutig
        referenzierbare Abschnitte gegliedert. Daraus werden quellengebundene
        Evidenz und gemeinsame fachliche Referenzgegenstände für mehrere
        Persona-Interpretationen abgeleitet. Übereinstimmungen, Abweichungen,
        mögliche Mehrdeutigkeiten und Informationslücken werden anschließend
        für die menschliche Prüfung zusammengeführt. Der LLM-Service
        unterstützt einzelne Interpretationsschritte, trifft jedoch keine
        fachliche Autorisierungsentscheidung. Die englischen Bezeichnungen
        in der Abbildung entsprechen den Element- und Funktionsnamen des
        SysML-v2-Modells. Eigene Darstellung auf Basis des entwickelten
        SysML-v2-Architekturmodells.}
    \label{fig:implementation_information_processing}
\end{figure}

Zu Beginn wird der Inhalt der registrierten Engineering-Quelle in eine
einheitliche interne Darstellung überführt. Dieser reproduzierbare
Aufbereitungsschritt wird im Modell als \emph{Source Projection} bezeichnet.
Er verändert die fachliche Aussage der Quelle nicht und fügt keine neue
fachliche Interpretation hinzu. Gleichzeitig bleibt festgehalten, aus
welcher Quelle und aus welchem konkreten Verarbeitungslauf die aufbereiteten
Informationen stammen.

Die aufbereitete Quelle wird anschließend in eindeutig referenzierbare
Abschnitte gegliedert, die im Modell als \emph{Source Analysis Units}
bezeichnet werden. Innerhalb dieser Abschnitte werden potenziell relevante
Engineering-Informationen identifiziert und weiterhin mit ihren konkreten
Fundstellen in der Ausgangsquelle verknüpft. Diese quellengebundenen
Beleginformationen entsprechen der in
Abschnitt~\ref{sec:architecture_grounded_interpretation} eingeführten
quellengebundenen Evidenz. In der Modellansicht wird hierfür die Bezeichnung
\emph{source-grounded Evidence} verwendet.

Aus der quellengebundenen Evidenz werden anschließend fachliche
Referenzgegenstände (\emph{Engineering Subjects}) gebildet. Sie sorgen dafür,
dass die nachfolgenden Interpretationen jeweils denselben fachlichen
Gegenstand betrachten und ihre Ergebnisse miteinander verglichen werden
können. Damit wird die in
Abschnitt~\ref{sec:architecture_grounded_interpretation} entwickelte
Trennung zwischen der Identifikation einer gemeinsamen quellengebundenen
Informationsgrundlage und ihrer multiperspektivischen Interpretation
umgesetzt.

Für die Interpretation werden drei konfigurierte Personas eingesetzt.
Sie betrachten jeweils denselben fachlichen Referenzgegenstand, setzen jedoch
unterschiedliche Schwerpunkte: Eine Persona orientiert sich möglichst eng
an der vorliegenden Quelle, eine zweite betrachtet den Inhalt aus
Systems-Engineering-Sicht und eine dritte prüft gezielt auf
Mehrdeutigkeiten und fehlende Informationen. Im Prototyp werden diese
Personas als \emph{Literal Interpreter},
\emph{Systems Engineering Interpreter} und
\emph{Skeptical Ambiguity Interpreter} bezeichnet. Die vollständigen
Persona-Konfigurationen sind in
Anhang~\ref{app:personas_interpretation} dokumentiert.

Die Persona-Interpretationen werden anschließend jeweils für denselben
fachlichen Referenzgegenstand miteinander verglichen. Dabei werden die in
Abschnitt~\ref{sec:architecture_grounded_interpretation} beschriebenen
Zustände Übereinstimmung (\emph{Consensus}), Varianz
(\emph{Variance}) und gemeinsame Mehrdeutigkeit
(\emph{Shared Ambiguity}) bestimmt. Sie zeigen, wo die Perspektiven zu
vergleichbaren Ergebnissen gelangen, voneinander abweichen oder dieselbe
Unsicherheit erkennen. Diese Informationen unterstützen die spätere
menschliche Prüfung, stellen selbst jedoch noch keine fachliche
Autorisierung dar.

Zusätzlich werden Informationen zur verwendeten Beleggrundlage und zu
möglichen Informationslücken ermittelt. Die im Modell als
\emph{Coverage} bezeichnete Information macht sichtbar, auf welche Bereiche
der Quelle sich die Verarbeitung stützt. \emph{Evidence Sufficiency}
beschreibt, ob für den betrachteten Gegenstand eine ausreichende
Beleggrundlage vorliegt oder relevante Informationslücken bestehen.
Semantische Referenzen und festgelegte Terminologie unterstützen darüber
hinaus eine konsistente Einordnung der Inhalte. Sie ergänzen die
quellengebundene Evidenz, ersetzen jedoch weder die ursprüngliche
Informationsgrundlage noch die nachfolgende menschliche Bewertung.

Die für die fachliche Prüfung relevanten Informationen werden anschließend
zu einem gemeinsamen Prüfzustand zusammengeführt, der im Modell als
\emph{Human Review Package} bezeichnet wird. Er enthält den betrachteten
fachlichen Referenzgegenstand, die zugehörige quellengebundene Evidenz, die
Persona-Interpretationen sowie Informationen zu Übereinstimmungen,
Abweichungen, Mehrdeutigkeiten und möglichen Informationslücken. Dieser
Zustand wird dauerhaft gespeichert und bildet die Grundlage für die
nachfolgende menschliche Prüfung.

Bis zu diesem Punkt wird jede Engineering-Quelle für sich verarbeitet.
Die Verarbeitung erzeugt damit eine nachvollziehbare Grundlage für die
fachliche Entscheidung, autorisiert die abgeleiteten Engineering-Informationen
aber noch nicht selbst. Die menschliche Autorisierung, die spätere
projektweite Zusammenführung mehrerer Quellen und die Ableitung konkreter
Modellelemente erfolgen in den nachgelagerten Verarbeitungsschritten.
Der folgende Abschnitt beschreibt zunächst die menschlichen Entscheidungen
über die fachliche Information und ihre Verwendung im Modell.

\subsection{Realisierung der menschlichen Autorisierungsgrenzen}
\label{sec:implementation_authority}

Die in Abschnitt~\ref{sec:implementation_information_processing}
beschriebene Verarbeitung bereitet Engineering-Information für eine
fachliche Prüfung auf, autorisiert die erzeugten Ergebnisse jedoch noch
nicht. Der Prototyp setzt deshalb das in
Abschnitt~\ref{sec:architecture_human_authority} entwickelte Konzept der
\emph{Human Authority} an mehreren Stellen des Verarbeitungsablaufs um.
Dabei entscheidet ein Mensch jeweils über einen konkret vorliegenden
Informations- oder Modellzustand, bevor dieser für den nächsten fachlich
relevanten Verarbeitungsschritt verwendet werden darf.

Die erste dieser Autorisierungsgrenzen betrifft die fachliche Verwendung
aufbereiteter Engineering-Information. Eine zweite betrifft deren
strukturelle Einordnung in das Modell. Eine weitere menschliche Entscheidung
erfolgt am Ende des Verarbeitungspfads vor der Freigabe des erzeugten
SysML-v2-Modells. Damit werden unterschiedliche Entscheidungsgegenstände
nicht in einem einzigen abschließenden Review zusammengefasst.

Im CATIA-Modell sind die dafür erforderlichen Funktionen im Subsystem
\emph{SBS\_05 Human Review and Engineering Authority} zusammengeführt.
Abbildung~\ref{fig:implementation_authority_rfl} zeigt einen Ausschnitt
dieser Modellierung und stellt die Verbindung zwischen ausgewählten
Anforderungen, den daraus abgeleiteten Funktionen und den zugehörigen
logischen Systemelementen dar.

\begin{figure}[H]
    \centering
    \includegraphics[
        width=\textwidth,
        height=0.82\textheight,
        keepaspectratio
    ]{Graphiken/5.5_Human_Authority_RFL.pdf}
    \caption[R/F/L-Zusammenhang der menschlichen Autorisierungsgrenzen]{
        Ausschnitt der R/F/L-Modellierung für die menschlichen
        Autorisierungsgrenzen. Dargestellt sind ausgewählte Anforderungen,
        Funktionen und logische Systemelemente für die fachliche
        Autorisierung von Engineering-Information, die menschlich
        autorisierte Modellplatzierung sowie die abschließende Modellprüfung
        und Freigabe. Die englischen Bezeichnungen entsprechen den Namen der
        Elemente im SysML-v2-Modell. Eigene Darstellung auf Basis des
        entwickelten SysML-v2-Architekturmodells.}
    \label{fig:implementation_authority_rfl}
\end{figure}

Die erste Autorisierungsgrenze liegt zwischen den in
Abschnitt~\ref{sec:implementation_information_processing} erzeugten
Bewertungsinformationen und der fachlich autorisierten
Engineering-Information. Für die menschliche Prüfung stellt der Prototyp
den zuvor aufbereiteten fachlichen Gegenstand gemeinsam mit seinen
Quellenbezügen, Persona-Interpretationen sowie den festgestellten
Übereinstimmungen, Abweichungen und Mehrdeutigkeiten bereit. Der konkret
geprüfte Inhalt wird dabei als eigener versionierter Zustand gespeichert.

Die menschliche Entscheidung wird genau an diesen geprüften Zustand
gebunden. Wird der entscheidungsrelevante Inhalt anschließend verändert,
ist deshalb eine erneute Prüfung und Autorisierung erforderlich. Eine
bestehende Entscheidung wird nicht automatisch auf einen veränderten
Inhalt übertragen. Damit wird die in
Abschnitt~\ref{sec:architecture_human_authority} entwickelte Bindung einer
menschlichen Entscheidung an den tatsächlich geprüften Zustand technisch
umgesetzt.

Erst durch diese Entscheidung entsteht fachlich autorisierte
Engineering-Information. Im Modell wird der entsprechende Zustand als
\emph{Approved Engineering Information} bezeichnet. Die zuvor erzeugten
Persona-Interpretationen sowie Informationen über Übereinstimmungen,
Varianz und Mehrdeutigkeiten bleiben dagegen Bestandteile der
Entscheidungsgrundlage. Sie unterstützen die menschliche Bewertung, besitzen
aber selbst keine fachliche Autorität. Die hierfür zuständige
Entscheidungsbefugnis entspricht der in Kapitel~\ref{sec:architecture}
eingeführten \emph{Human Engineering Authority}.

Mit der fachlichen Autorisierung ist noch nicht entschieden, an welcher
Stelle des vorgesehenen Modellrahmens die Information repräsentiert werden
soll. Für modellrelevante Inhalte erzeugt der Prototyp deshalb zunächst
Vorschläge für diese strukturelle Einordnung. Im Modell wird dieser Schritt
als \emph{Model Placement} bezeichnet.

Auch diese Vorschläge können durch mehrere Modellierungsperspektiven
unterstützt werden. Im Prototyp werden dafür die in
Anhang~\ref{app:personas_modeling} dokumentierten Personas
\emph{SysML Profile Modeler}, \emph{System Architecture Modeler} und
\emph{Conservative Modeling Reviewer} eingesetzt. Sie arbeiten auf
derselben autorisierten Engineering-Information und denselben zulässigen
Positionen innerhalb des Modellrahmens. Ihre Ergebnisse werden für die
menschliche Entscheidung zusammengeführt, autorisieren die
Modellplatzierung jedoch nicht selbst.

Die strukturelle Einordnung wird erst durch eine explizite menschliche
Modellentscheidung autorisiert. Diese Entscheidung wird an den betreffenden
Engineering-Inhalt sowie an die verwendete Version der zulässigen
Modellstruktur gebunden. Diese Modellstruktur ist im Prototyp in einem
versionierten \emph{Model Structure Profile} beschrieben. Die daraus
entstehenden autorisierten Platzierungsentscheidungen bilden anschließend
die Grundlage für den Aufbau des internen Modells. Diese
Entscheidungsbefugnis entspricht der in
Abschnitt~\ref{sec:architecture_human_authority} eingeführten
\emph{Model Authority}. Die nachgelagerte Modellbildung wird in
Abschnitt~\ref{sec:implementation_model_realization} beschrieben.

Eine weitere menschliche Entscheidungsgrenze liegt am Ende der
Modellerzeugung. Dort wird das erzeugte SysML-v2-Artefakt gemeinsam mit dem
Ergebnis seiner Validierung betrachtet. Die Validierung prüft die dafür
festgelegten formalen und technischen Bedingungen, gibt das Modell aber
nicht selbst frei. Erst die anschließende menschliche Modellprüfung und
Freigabe autorisiert den konkret geprüften Modellzustand für die
Veröffentlichung.

Die technische Ablaufsteuerung prüft anschließend lediglich, ob die für
eine Veröffentlichung erforderlichen Zustände und Entscheidungen vorliegen.
Sie trifft selbst keine fachliche Freigabeentscheidung.

Die prototypische Realisierung unterscheidet damit drei menschlich
autorisierte Übergänge: die fachliche Autorisierung von
Engineering-Information, die Autorisierung ihrer strukturellen Einordnung
in das Modell sowie die abschließende Prüfung und Freigabe des erzeugten
Modellzustands. Jede Entscheidung bleibt mit dem jeweils geprüften Zustand
verknüpft. Wie diese Zustände und ihre Beziehungen dauerhaft gespeichert
und nachvollziehbar gehalten werden, beschreibt
Abschnitt~\ref{sec:implementation_persistence_limits}.

\subsection{Modellbildung und SysML-v2-Realisierung}
\label{sec:implementation_model_realization}

Nach der fachlichen Autorisierung der Engineering-Information und der
Entscheidung über ihre strukturelle Einordnung kann der eigentliche
Modellaufbau erfolgen. Die nachfolgenden Schritte verändern die zuvor
autorisierten fachlichen und modellbezogenen Entscheidungen nicht mehr,
sondern führen sie nach festgelegten Regeln zu einem internen Modell und
anschließend zu einer SysML-v2-Repräsentation zusammen.

Aus der autorisierten Engineering-Information werden zunächst
\emph{Model Candidates} abgeleitet. Wie in
Abschnitt~\ref{sec:architecture_meaning_representation} beschrieben,
repräsentieren sie mögliche Modellelemente und Beziehungen. Für den
Modellaufbau werden ausschließlich diejenigen Model Candidates verwendet,
deren strukturelle Einordnung zuvor durch die Model Authority autorisiert
wurde. Nicht autorisierte oder zurückgestellte Vorschläge werden nicht in
das Modell übernommen.

Die autorisierten Modellinhalte werden anschließend zu einem
\emph{Internal Engineering Model} (IEM) zusammengeführt. Das IEM bildet,
wie in Abschnitt~\ref{sec:architecture_meaning_representation} eingeführt,
eine interne Modellrepräsentation zwischen den autorisierten
Engineering-Informationen und ihrer späteren Darstellung in SysML~v2.
Elemente und Beziehungen erhalten darin stabile interne Identitäten und
bleiben mit den fachlichen Informationen und Modellentscheidungen
verknüpft, aus denen sie hervorgegangen sind.

Abbildung~\ref{fig:implementation_model_realization} zeigt die wesentlichen
Übergänge von der autorisierten Modellplatzierung über den Aufbau des IEM
bis zur SysML-v2-Erzeugung und anschließenden Validierung.

\begin{figure}[H]
    \centering
    \includegraphics[
        width=\textwidth,
        height=0.82\textheight,
        keepaspectratio
    ]{Graphiken/5.6_Model_Realization_Cross_Subsystem.pdf}
    \caption[Modellbildung und SysML-v2-Realisierung]{
        Modellbildung und SysML-v2-Realisierung über mehrere Subsysteme.
        Aus fachlich autorisierter Engineering-Information werden zunächst
        Vorschläge für Modellelemente und deren strukturelle Einordnung
        abgeleitet. Nach der menschlichen Autorisierung der Modellplatzierung
        werden die freigegebenen Inhalte im Internal Engineering Model
        zusammengeführt und anschließend regelgebunden in SysML~v2
        überführt. Das erzeugte Artefakt wird getrennt validiert und danach
        der abschließenden menschlichen Modellprüfung und Freigabe
        bereitgestellt. Die englischen Bezeichnungen entsprechen den
        Element- und Funktionsnamen des SysML-v2-Modells.
        Eigene Darstellung auf Basis des entwickelten
        SysML-v2-Architekturmodells.}
    \label{fig:implementation_model_realization}
\end{figure}

Der Aufbau des IEM erfolgt nach festgelegten Regeln. Als strukturelles
Grundgerüst dient eine versionierte Vorlage, die im Prototyp als
\emph{Framework Template} bezeichnet wird. Sie legt die vorgesehene
Grundstruktur des Modells fest. Die autorisierten Model Candidates werden
entsprechend ihrer zuvor freigegebenen Platzierungsentscheidungen in diese
Struktur eingeordnet. Dabei wird ihre fachliche Bedeutung nicht erneut
interpretiert.

Kann ein autorisierter Model Candidate nicht eindeutig entsprechend der
vorliegenden Entscheidung eingeordnet werden, wird der Modellaufbau an
dieser Stelle unterbrochen. Die offene Frage muss dann durch eine erneute
menschliche Modellentscheidung geklärt werden. Der Prototyp ergänzt in
diesem Fall nicht selbstständig zusätzliche Hierarchien, Beziehungen oder
strukturelle Zuordnungen.

Jede inhaltlich unterschiedliche autorisierte Fassung des IEM wird als
eigener unveränderlicher Stand gespeichert. Dadurch bleibt nachvollziehbar,
welche Model Candidates, Platzierungsentscheidungen, Modellstruktur und
Regeln zu einer konkreten Fassung des internen Modells geführt haben.

Das IEM ist noch keine SysML-v2-Repräsentation. Es enthält die autorisierten
Modellelemente und Beziehungen zunächst unabhängig von ihrer konkreten
textuellen Darstellung in SysML~v2. Damit bleibt die bereits in
Abschnitt~\ref{sec:architecture_meaning_representation} entwickelte
Trennung zwischen internem Modellinhalt und zielmodellspezifischer
Formulierung erhalten.

Vor der eigentlichen SysML-v2-Erzeugung kann die Formulierung des
Modellinhalts hinsichtlich seiner Darstellungsqualität überprüft werden.
Dafür wird ein versionierter Satz von Qualitätskriterien verwendet. Die in
Anhang~\ref{app:persona_model_quality} dokumentierte Persona
\emph{Conservative Engineering Model Editor} kann auf dieser Grundlage
einen LLM-gestützten Verbesserungsvorschlag erzeugen. Dabei dürfen
beispielsweise Bezeichnungen oder Formulierungen präzisiert werden, jedoch
keine zusätzlichen Engineering-Informationen entstehen.

Diese Vorschläge werden anschließend menschlich geprüft und können
akzeptiert, verändert oder verworfen werden. Nur ein entsprechend
autorisierter Stand wird in das IEM übernommen. Dadurch bleibt die
zielmodellspezifische Formulierung von der vorausgehenden fachlichen
Autorisierung und Modellplatzierung getrennt.

Für die anschließende SysML-v2-Erzeugung verwendet der Prototyp drei
versionierte Regelwerke, die bereits in
Abschnitt~\ref{sec:architecture_controlled_transformation} eingeführt
wurden. Das \emph{Target Notation Profile} legt fest, welche
SysML-v2-Konstrukte vom Prototyp unterstützt werden. Das
\emph{Generation Profile} beschreibt, wie die internen Modellkonzepte auf
diese Konstrukte abgebildet werden. Das
\emph{Artifact Structure Profile} legt fest, wie die erzeugten Inhalte
beispielsweise in Packages und Dateien organisiert werden.

Vor der Erzeugung des SysML-v2-Artefakts prüft der Prototyp, ob für alle
zu überführenden Modellelemente und Beziehungen eine eindeutige und
unterstützte Abbildungsregel vorhanden ist. Fehlt eine solche Regel, wird
die Erzeugung unterbrochen. Der Prototyp ersetzt nicht unterstützte
Modellinhalte damit nicht automatisch durch andere SysML-v2-Konstrukte.

Auf Grundlage des IEM und der festgelegten Transformationsregeln wird
anschließend das SysML-v2-Artefakt erzeugt. Zusammen mit dem eigentlichen
SysML-v2-Text werden die Informationen gespeichert, die eine
Rückverfolgung auf das verwendete IEM und die angewendeten
Transformationsregeln ermöglichen.

Die Validierung des erzeugten SysML-v2-Artefakts erfolgt getrennt von seiner
Erzeugung. Zunächst führt der Prototyp interne regelbasierte Prüfungen durch.
Zusätzlich wird das erzeugte Artefakt mit der \emph{SYSIDE Modeler CLI} in
einer externen SysML-v2-Toolumgebung geprüft. Die konkrete Durchführung
und Bewertung dieser Prüfungen wird in
Kapitel~\ref{sec:validation} beschrieben.

Eine erfolgreiche Validierung stellt noch keine abschließende Freigabe des
Modells dar. Das erzeugte und validierte SysML-v2-Artefakt wird anschließend
der in Abschnitt~\ref{sec:implementation_authority} beschriebenen
menschlichen Modellprüfung und Freigabe vorgelegt.

Die prototypische Realisierung hält damit die Autorisierung der
Modellplatzierung, den Aufbau des IEM, die zielmodellspezifische
Formulierung, die SysML-v2-Erzeugung und die Validierung als voneinander
getrennte Verarbeitungsschritte aufrecht.

\subsection{Verarbeitung mehrerer Quellen und Modellevolution}
\label{sec:implementation_multisource}

Die bisher beschriebenen Verarbeitungsschritte beziehen sich zunächst auf
einzelne Engineering-Quellen. Für die gemeinsame Modellbildung müssen jedoch
auch Informationen aus mehreren Quellen projektweit berücksichtigt werden
können. Der Prototyp realisiert dafür die in
Abschnitt~\ref{sec:architecture_multisource} entwickelte Abfolge aus
quellenbezogener Verarbeitung, Projektzulassung, projektweitem semantischem
Abgleich und menschlicher Projektentscheidung.

Auch bei der gemeinsamen Verarbeitung bleibt nachvollziehbar, aus welcher
Engineering-Quelle eine Information stammt und in welchem Verarbeitungslauf
sie erzeugt wurde. Informationen verschiedener Quellen werden daher nicht
bereits während der quellenbezogenen Verarbeitung zusammengeführt. Erst
fachlich autorisierte Engineering-Information wird für die nachfolgenden
projektweiten Verarbeitungsschritte berücksichtigt. Die technische
Speicherung der hierfür erforderlichen Herkunfts- und
Verarbeitungsinformationen wird in
Abschnitt~\ref{sec:implementation_persistence_limits} beschrieben.

Vor der projektweiten Zusammenführung wird zunächst geprüft, ob eine
Engineering-Quelle für das betrachtete Projekt berücksichtigt werden soll.
Diese in Abschnitt~\ref{sec:architecture_multisource} als
\emph{Projektzulassung} (\emph{Project Fit}) eingeführte Prüfung verwendet
Informationen über das Projekt, die registrierten Quellen und bereits
autorisierte Engineering-Information. Das Ergebnis unterscheidet zwischen
einer plausiblen Projektzugehörigkeit, einer unklaren Zuordnung und einer
wahrscheinlich fehlenden Projektzugehörigkeit. Bei einem unklaren oder
wahrscheinlich nicht passenden Ergebnis ist eine menschliche Entscheidung
erforderlich, bevor die Quelle projektweit weiterverarbeitet werden darf.

Nach der Projektzulassung werden die fachlich autorisierten Informationen
unterschiedlicher Quellen miteinander verglichen. Als gemeinsame
Bezugspunkte dienen die bereits in
Abschnitt~\ref{sec:implementation_information_processing} beschriebenen
fachlichen Referenzgegenstände (\emph{Engineering Subjects}). Der Prototyp
gruppiert solche Referenzgegenstände, die sich möglicherweise auf denselben
fachlichen Zusammenhang beziehen. Anschließend wird automatisch geprüft, ob
jeder bereitgestellte Referenzgegenstand genau einmal berücksichtigt wurde
und keine unbekannten Referenzen hinzugekommen sind.

Für die daraus entstehenden Vergleichsfälle wird untersucht, in welcher
fachlichen Beziehung die beteiligten Informationen zueinander stehen.
Entsprechend Abschnitt~\ref{sec:architecture_multisource} unterscheidet der
Prototyp zwischen gleichbedeutenden (\emph{equivalent}), sich ergänzenden
(\emph{complementary}), möglicherweise widersprüchlichen
(\emph{potential conflict}), voneinander unabhängigen (\emph{distinct}) und
nicht eindeutig einzuordnenden (\emph{uncertain}) Beziehungen. Liegt für
einen fachlichen Referenzgegenstand kein entsprechender Beitrag aus einer
anderen Quelle vor, bleibt er als eigenständiger Fall erhalten.

Diese maschinell unterstützte Einordnung entscheidet noch nicht darüber,
welche Information projektweit verwendet werden soll. Sie bereitet vielmehr
die menschliche Projektentscheidung vor. Dabei kann entschieden werden, dass
Informationen unabhängig voneinander bestehen bleiben, gemeinsam gültig
sind, einander ablösen oder zunächst ungeklärt bleiben. Die bereits
quellenbezogen autorisierten Engineering-Informationen werden durch diese
Entscheidung nicht überschrieben.

Entsteht durch die projektweite Zusammenführung eine neue oder wesentlich
veränderte Engineering-Aussage, muss auch diese erneut fachlich geprüft und
autorisiert werden. Aus den anschließend für das Projekt gültigen
Engineering-Informationen entsteht ein autorisierter projektweiter
Engineering-Zustand. Im Modell wird dieser als
\emph{Project Engineering State} bezeichnet. Er bildet die fachliche
Grundlage für die weitere Modellbildung.

Aus diesem projektweiten Engineering-Zustand werden die für die
Modellbildung erforderlichen \emph{Model Candidates} abgeleitet. Besteht
bereits ein autorisierter Modellzustand, wird zusätzlich geprüft, welche
Auswirkungen die neuen oder geänderten Engineering-Informationen auf dieses
Modell haben. Dieser in
Abschnitt~\ref{sec:architecture_multisource} als
\emph{Modellwirkungsabgleich} beschriebene Schritt wird im Prototyp als
\emph{Model Impact Reconciliation} bezeichnet.

Der Modellwirkungsabgleich erzeugt Vorschläge für erforderliche
Modelländerungen, autorisiert diese jedoch nicht selbst. Erst eine
anschließende menschliche Modellentscheidung kann einen neuen
Modellzustand autorisieren. Der bisherige autorisierte Zustand bleibt dabei
erhalten. Auf diese Weise kann nachvollzogen werden, welche neuen oder
geänderten Engineering-Informationen zu einer Weiterentwicklung des Modells
geführt haben.

\subsection{Persistenz, Provenienz und Grenzen}
\label{sec:implementation_persistence_limits}

Die prototypische Realisierung speichert die für den Transformationsprozess
relevanten Informations-, Entscheidungs- und Modellstände dauerhaft. Diese
dauerhafte Speicherung wird im Folgenden als \emph{Persistenz} bezeichnet.
Gespeicherte Zustände werden nicht durch spätere Bearbeitungen überschrieben,
sondern bleiben als frühere Stände erhalten. Dadurch kann nachvollzogen
werden, aus welchen Quellen, Verarbeitungsschritten und menschlichen
Entscheidungen ein späteres Modell entstanden ist.

Bereits bei der Verarbeitung einer Engineering-Quelle wird deshalb
festgehalten, welche konkrete Quelle verarbeitet wurde, zu welchem Projekt
sie gehört und unter welcher Verarbeitungskonfiguration der jeweilige Lauf
durchgeführt wurde. Jeder Verarbeitungslauf erhält eine eindeutige Identität.
Technische Wiederholungen innerhalb desselben Laufs werden separat
gespeichert. Ändert sich dagegen die zugrunde liegende Quelle oder eine für
die Verarbeitung wesentliche Konfiguration, entsteht ein neuer
Verarbeitungslauf.

Für die technische Umsetzung verwendet der Prototyp hierfür unter anderem
Manifeste, Referenzen und inhaltsgebundene Prüfwerte
(\emph{Fingerprints}). Diese ermöglichen es, gespeicherte Zustände eindeutig
miteinander zu verknüpfen und spätere Änderungen festzustellen.

Dieses Prinzip wird über die weiteren Verarbeitungsschritte fortgeführt.
Auch menschliche Prüfentscheidungen, autorisierte Engineering-Information,
Modellentscheidungen, interne Modellstände, erzeugte SysML-v2-Artefakte und
Validierungsergebnisse bleiben jeweils mit den Zuständen verknüpft, aus denen
sie hervorgegangen sind. Tabelle~\ref{tab:implementation_persisted_artifacts}
fasst die wichtigsten gespeicherten Informationen zusammen.

\begin{table}[htbp]
    \centering
    \small
    \caption[Persistierte Informationen des Prototyps]{
        Zentrale persistierte Informationen und ihre Funktion}
    \label{tab:implementation_persisted_artifacts}
    \begin{tabular}{p{3.2cm}p{5.0cm}p{5.7cm}}
        \hline
        \textbf{Bereich} &
        \textbf{Gespeicherte Informationen} &
        \textbf{Funktion} \\
        \hline

        Quelle und Verarbeitung &
        Quelle, aufbereitete Quellenrepräsentation,
        Verarbeitungslauf, technische Wiederholungen und Verlauf &
        Nachvollziehbare Zuordnung von Eingangsinformation,
        Verarbeitungskonfiguration und Verarbeitungshistorie \\

        Interpretation &
        quellengebundene Evidenz, Engineering Subjects,
        Persona-Interpretationen, Übereinstimmungen,
        Abweichungen und Mehrdeutigkeiten &
        Nachvollziehbare Grundlage für die anschließende
        menschliche Prüfung \\

        Fachliche Autorisierung &
        geprüfter Informationsstand und menschliche
        Autorisierungsentscheidung &
        Bindung der fachlichen Freigabe an genau den Zustand,
        der tatsächlich geprüft wurde \\

        Projektweite Verarbeitung &
        Projektzulassung, quellenübergreifende Vergleichsfälle,
        menschliche Projektentscheidungen und projektweiter
        Engineering-Zustand &
        Nachvollziehbare Zusammenführung mehrerer Quellen unter
        Erhalt ihrer jeweiligen Herkunft \\

        Modellbildung &
        Model Candidates, Platzierungsentscheidungen,
        Modellwirkungsabgleich und Internal Engineering Model &
        Nachvollziehbarer Aufbau und Weiterentwicklung des
        internen Modells \\

        SysML-v2-Erzeugung und Freigabe &
        erzeugtes SysML-v2-Artefakt, Validierungsergebnis
        und abschließende Freigabeentscheidung &
        Rückverfolgbarkeit vom freigegebenen Artefakt bis zu
        Modell-, Engineering- und Quellenzuständen \\
        \hline
    \end{tabular}
\end{table}

Die Provenienz ergibt sich aus diesen fortlaufenden Verknüpfungen. Ein
erzeugtes SysML-v2-Artefakt kann dadurch auf das verwendete IEM, die
zugrunde liegenden Modell- und Engineering-Entscheidungen sowie schließlich
auf die beteiligten Engineering-Quellen zurückgeführt werden. Gleichzeitig
bleibt sichtbar, welche Zwischenstände und menschlichen Entscheidungen an
seiner Entstehung beteiligt waren.

\paragraph{Technische Grenzen des Prototyps}

Die prototypische Realisierung ist auf einen bewusst abgegrenzten
Funktionsumfang beschränkt. Tabelle~\ref{tab:implementation_limits} fasst
die für die weitere Bewertung wesentlichen technischen Grenzen zusammen.

\begin{table}[htbp]
    \centering
    \small
    \caption[Technische Grenzen des Prototyps]{
        Wesentliche technische Grenzen der prototypischen Realisierung}
    \label{tab:implementation_limits}
    \begin{tabular}{p{3.3cm}p{4.2cm}p{6.2cm}}
        \hline
        \textbf{Bereich} &
        \textbf{Aktueller Stand} &
        \textbf{Auswirkung} \\
        \hline

        Verarbeitung mehrerer Quellen &
        Mehrere Quellen können innerhalb eines gemeinsamen
        Projektkontexts verarbeitet und zusammengeführt werden &
        Neu hinzukommende Informationen führen nicht automatisch
        zu einer vollständigen Neubewertung sämtlicher bereits
        vorhandener Projektinformationen \\

        Eingabeformate &
        Textorientierte Engineering-Information &
        Unterstützt werden Text, Markdown, JSON, CSV/TSV,
        deterministische Textrepräsentationen und PDF-Textlayer.
        OCR, rein bildbasierte Dokumente, Zeichnungsinterpretation
        und allgemeine multimodale Verarbeitung sind nicht
        Bestandteil des Prototyps \\

        Modellierungsrahmen &
        Vorgegebene RFLP-Struktur &
        Modellableitung und strukturelle Einordnung erfolgen
        innerhalb des für den Prototyp definierten
        Stakeholder-, System- und Subsystemrahmens. Beliebige
        Zielstrukturen werden nicht unterstützt \\

        SysML-v2-Umfang &
        Begrenzter unterstützter Ausschnitt der Zielsprache &
        Nur ausdrücklich unterstützte SysML-v2-Konstrukte und
        Beziehungen können regelgebunden erzeugt werden \\

        Externe Validierung &
        SYSIDE Modeler CLI erforderlich &
        Für die vollständige externe Prüfung muss eine kompatible
        SYSIDE-Toolumgebung verfügbar sein \\

        Menschliche Autorisierung &
        Bestandteil des vorgesehenen Verarbeitungsablaufs &
        Fachlich oder modellbezogen relevante Entscheidungen sowie
        die abschließende Freigabe erfordern explizite menschliche
        Entscheidungen; der Prototyp arbeitet daher nicht als
        autonomer Modellgenerator \\
        \hline
    \end{tabular}
\end{table}

Diese Grenzen bestimmen den technischen Untersuchungsrahmen der
prototypischen Realisierung. Ihre Auswirkungen auf die Aussagekraft,
Übertragbarkeit und Generalisierbarkeit der Ergebnisse werden in
Kapitel~\ref{chap:discussion} diskutiert.